\documentclass[12pt]{article}
\usepackage[utf8]{inputenc}
\usepackage{graphicx}
\usepackage{hyperref}
\usepackage{amssymb}
\usepackage{amsmath}
\usepackage[T1]{fontenc}
\usepackage{lmodern}
\usepackage{textcomp}
\usepackage{pgfplots}
\usepackage{pgfplotstable}
\usepackage{longtable}
\usepackage[capitalise]{cleveref}
\usepackage{caption}
\pgfplotsset{compat=1.18}
\captionsetup[table]{position=top,justification=raggedright,singlelinecheck=false}
\captionsetup[longtable]{position=top,justification=raggedright,singlelinecheck=false}
\captionsetup[figure]{justification=raggedright,singlelinecheck=false}
\setlength\LTleft{0pt}
\setlength\LTright{0pt}
\setlength\LTcapwidth{\textwidth}
\pgfplotstableset{
  show variable once/.style={
    columns/variable/.style={
      string type,
      column name=Variable,
      postproc cell content/.code={
        \pgfkeysgetvalue{/pgfplots/table/@cell content}{\cellcontent}
        \edef\currentvar{\cellcontent}
        \ifx\currentvar\lastvar
          \pgfkeyssetvalue{/pgfplots/table/@cell content}{}
        \else
          \ifnum\pgfplotstablerow>0
            \pgfkeyssetevalue{/pgfplots/table/@cell content}{\noexpand\hline\space\cellcontent}
          \fi
          \xdef\lastvar{\currentvar}
        \fi
      }
    }
  }
}
\def\lastvar{}
\setlength{\parskip}{0.8em}
\setlength{\parindent}{0pt}
\newcommand{\indep}{\perp \!\!\! \perp}%\renewcommand{\thesubsection}{\thesection.\alph{subsection}}

\title{\Large\bfseries ASSIGNMENT 5: LINEAR MIXED EFFECTS MODEL}
\author{Chenghao Wang}
\date{\today}

\begin{document}

\maketitle

\section{Model Specification and Fitting}

\subsection{Model formulation}

The model formulation is as follows:

\begin{multline}
\label{eq:final-model}
Y_{i} = \beta_{0} + I(TG_i = bud)\,visit_{i}\beta_{visit \times bud}
+ I(TG_i = ned)\,visit_{i}\beta_{visit \times ned} \\
{}+ I(TG_i = plbo)\,visit_{i}\beta_{visit \times plbo}
+ age\_rz_{i}\,\beta_{age\_rz}
+ I(gender_{i} = m)\,\beta_{gender}\\
{}+ I(ethnic_{i} = b)\,\beta_{ethnic:black}
+ I(ethnic_{i} = h)\,\beta_{ethnic:hispanic}
{}+ I(ethnic_{i} = o)\,\beta_{ethnic:other} \\
+ b_{0,i} + visit_i b_{1,i} + \varepsilon_{i}\\
(b_{0,i}, b_{1,i})^T \equiv b_i \sim N(0, D),\,
\varepsilon_{i} \sim N(0,\sigma^2_{TG_i}I).
\end{multline}

The fixed effects include intercept (at baseline), visit time
(three group-specific slopes), age at randomization, gender,
and ethnicity. The random effects include subject-specific random intercept
and subject-specific random slope for visit time.

We assume that $b_i \indep \varepsilon_i$.
In addition, we assume that the covariance matrix of $b_i$
is an unstructured matrix, and the covariance matrix of
$\varepsilon_i$ is a diagonal matrix with group-specific variances.

\subsection{Random effects structure}

A random effects structure with both random intercept and
random slope for visit time is chosen,
to account for the subject-level variability in
both the baseline outcome and the rate of change over visit time,
as we assume that the visit time effect is linear
in the fixed effects structure.
This is also justified because the measurements
were taken at multiple time points for each subject,
and thus it is reasonable to expect subject-specific trajectories.

R package \texttt{nlme} was used.
The model was fitted using REML.
The fixed effect estimates are demonstrated in \cref{tab:effect-table}.
The variance component estimates are shown in \cref{tab:random-effect-covariance}
and \cref{tab:residual-variance-group}.



\begin{table}[htbp]
\caption{Estimated fixed effects from the fitted mixed model.}
\centering
\pgfplotstabletypeset[
  col sep=comma,
  string replace*={"}{},
  display columns/0/.style={string type, column name=Effect, column type=l},
  display columns/1/.style={column name=Estimate, column type=r, fixed, precision=3},
  display columns/2/.style={column name=SE, column type=r, fixed, precision=4},
  display columns/3/.style={column name=Test statistic, column type=r},
  display columns/4/.style={string type, column name=\textit{p}-value, column type=r},
  every head row/.style={before row=\hline, after row=\hline},
  every last row/.style={after row=\hline}
]{../hw5_code/output/effect_table.csv}
\label{tab:effect-table}
\end{table}

\begin{table}[htbp]
\caption{Estimated random-effects covariance matrix.}
\centering
\pgfplotstabletypeset[
  col sep=comma,
  string replace*={"}{},
  display columns/0/.style={string type, column name=Random effect, column type=l},
  display columns/1/.style={column name=(Intercept), column type=r, fixed, precision=6},
  display columns/2/.style={column name=visitc, column type=r, fixed, precision=6},
  every head row/.style={before row=\hline, after row=\hline},
  every last row/.style={after row=\hline}
]{../hw5_code/output/random_effect_covariance_matrix.csv}
\label{tab:random-effect-covariance}
\end{table}

\begin{table}[htbp]
\caption{Estimated residual variance by treatment group.}
\centering
\pgfplotstabletypeset[
  col sep=comma,
  string replace*={"}{},
  display columns/0/.style={string type, column name=Treatment group, column type=l},
  display columns/1/.style={column name=Residual variance, column type=r, fixed, precision=5},
  every head row/.style={before row=\hline, after row=\hline},
  every last row/.style={after row=\hline}
]{../hw5_code/output/residual_variance_by_group.csv}
\label{tab:residual-variance-group}
\end{table}



\section{Evaluating the Random Effects}

\subsection{Testing random effects}

First, we evaluated the need for the random slope for visit time.
We compared the full model with a reduced model without
the random slope for visit time.
The likelihood ratio test (LRT) statistic is 3691.1,
and the p-value is < 0.001 (The null distribution
is $.5 \chi_{1}^2 + .5 \chi_{2}^2$.),
which indicates that the random slope for visit time is needed.
Next, we evaluated the need for the random intercept
by comparing the full model with a reduced model
without the random intercept. The LRT statistic is
10092.1, and the p-value is < 0.001,
which suggest that the random intercept is also needed.

We also evaluated the need for the random effects
using AIC.
As shown in \cref{tab:aic-comparison-random-effects},
the model with both random intercept and random slope for visit time
has the lowest AIC, which also supports the choice of the random effects
structure.

\begin{table}[htbp]
\caption{AIC comparison for random-effects structures.}
\centering
\pgfplotstabletypeset[
  col sep=comma,
  string replace*={"}{},
  display columns/0/.style={string type, column name=Model, column type=l},
  display columns/1/.style={column name=DF, column type=r},
  display columns/2/.style={column name=AIC, column type=r, fixed, precision=1},
  every head row/.style={before row=\hline, after row=\hline},
  every last row/.style={after row=\hline}
]{../hw5_code/output/aic_comparison.csv}
\label{tab:aic-comparison-random-effects}
\end{table}




\subsection{Variance components interpretation}

Under the model we specified,
if we express $D$ as
\begin{equation}
D = \begin{bmatrix}
\tau_{00} & \tau_{01} \\
\tau_{10} & \tau_{11}
\end{bmatrix},
\end{equation}

then the $(j, k)$ element of $Cov(Y_i)$ is as follows:

\begin{equation}
\tau_{00} + (visit_{ij} + visit_{ik})\tau_{01} + visit_{ij} visit_{ik}
\tau_{11} + I(j = k)\sigma^2_{TG_i}.
\end{equation}

This intuitively shows the contribution of the variance components.
Since the estimate of $\tau_{01}$ is positive,
we can tell that the variance of $Y_{ij}$
increases as visit time increases;
the covariance between $Y_{ij}$ and $Y_{ik}$ also increases
as visit time increases, which is consistent with
our observation in the previous Assignments.
Nonetheless, the characteristic that
the covariance between two measurements decreases as the
lag increases is not correctly modeled by this model.

The estimated variance components ($\hat{D}$) tells
us about the between-subject variability in the data.

\section{Subject-Specific Predictions}
\subsection{EBLUPs}

The empirical BLUPs (EBLUPs) for each subject
are displayed in \cref{fig:random-intercept-histogram}
and \cref{fig:random-slope-histogram}.

\begin{figure}[htbp]
\centering
\includegraphics[width=0.75\textwidth]{../hw5_code/output/random_intercept_histogram.png}
\caption{Histogram of EBLUPs for the random intercept.}
\label{fig:random-intercept-histogram}
\end{figure}

\begin{figure}[htbp]
\centering
\includegraphics[width=0.75\textwidth]{../hw5_code/output/random_slope_histogram.png}
\caption{Histogram of EBLUPs for the random slope of visit time.}
\label{fig:random-slope-histogram}
\end{figure}


\subsection{Subject-specific trajectories}

The subject-specific observed and fitted trajectories 
are illustrated in \cref{fig:subject-specific-trajectories},
for a subset of subjects. The fitted trajectories
were obtained by using both fixed effect estimates and EBLUPs.
The model captures individual patterns reasonably well.

\begin{figure}[htbp]
\centering
\includegraphics[width=0.9\textwidth]{../hw5_code/output/subject_specific_trajectory_plot.png}
\caption{Subject-specific predicted trajectories from the fitted mixed model.}
\label{fig:subject-specific-trajectories}
\end{figure}


\section{Residual and Influence Diagnostics}
\subsection{Residual analysis}

The R-P plot made using \texttt{nlme} package
is shown as \cref{fig:fit-diagnostic-plot}.
Although the vertical axis label shows "Standardized residuals",
the residuals plotted are actually the Pearson residuals
(The Pearson residuals are defined as 
$(y_{ij} - \hat{y}_{ij})/\sqrt{\widehat{Var}(y_{ij})}$,
while the standardized residuals are defined as
$(y_{ij} - \hat{y}_{ij})/\sqrt{\widehat{Var}(y_{ij} - \hat{y}_{ij})}$,
where $\hat{y}_{ij}$ is estimated using both fixed effect estimates
and EBLUPs). We find that the distribution of the Pearson residuals is approximately
symmetric. However, the dispersion of them increases as the fitted value
increases. The fitted values were also computed using both
fixed effect estimates and EBLUPs. Note that the fitted value
increases as visit time increases, and in Assignment 3, we showed
that the variance increases over time until month 60, and then tends to
be constant. Thus, we conclude that our model does not perfectly
capture the time-variance trend.

\begin{figure}[htbp]
\centering
\includegraphics[width=0.9\textwidth]{../hw5_code/output/fit_diagnostic_plot.png}
\caption{Residual diagnostic plots for the fitted linear mixed model.}
\label{fig:fit-diagnostic-plot}
\end{figure}


\subsection{Influence diagnostics}

The subject-level Cook's distance is illustrated
in \cref{fig:subject-level-cooks-distance}, where
the subjects are ordered by Cook's distance.
When computing the Cook's distance,
models were refitted with a limit of the number of
maximum iterations being set, for computational efficiency.
We find that there are two subjects (id 803 and id 345) with markedly larger
Cook's distance than the others.
These two subjects warrant further investigation.

\begin{figure}[htbp]
\centering
\includegraphics[width=0.9\textwidth]{../hw5_code/output/subject_level_cooks_distance.png}
\caption{Subject-level Cook's distance for the fitted linear mixed model.}
\label{fig:subject-level-cooks-distance}
\end{figure}


\section{Summary and Interpretation}

We fit a linear mixed effects model, including
a random intercept and a random slope for visit time.
The model's fixed effect structure includes
an intercept and three group-specific slopes for visit time,
and it adjusts for age at randomization, gender,
and ethnicity.

We performed a test of parallelism,
with $H_0: \beta_{visit \times bud} = \beta_{visit \times ned} = \beta_{visit \times plbo}$
versus $H_1: \text{at least one of these equalities fails}$.
The Wald $\chi^2$ statistic is 5.82,
and the p-value is 0.071, which indicates an insignificant
result. Thus we do not reject the null and conclude
that mean PostBD FEV1 does not change differently
over visit time among the three treatment groups.

We have shown that this model does not
perfectly capture the time-variance trend by R-P plot.
We tried to fit a model with heterogeneous residual variance
across visit time, but the model did not converge.
Besides, this model also does not
capture the characteristic that the covariance between two measurements
decreases as the lag increases.
Moreover, we find two influential subjects
that warrant further investigation.

\section{Code Availability}

The code is available at
\url{https://github.com/CHENGHAO-WANG/767_homework/tree/main/hw5_code}.


\end{document}
